\documentclass[UTF8, 12pt, fontset=fandol, oneside]{ctexart}
\usepackage{researchnotes}

\graphicspath{{tex/强化学习/figures/math-15/}}

\title{强化学习：心理学}
\author{}
\date{}

\begin{document}

\maketitle

\section*{第三部分\quad 深入考察}

在本书最后一部分中，我们越过前两部分介绍的标准强化学习思想，简要考察它们与心理学和神经科学的关系，抽样介绍若干强化学习应用，并讨论未来强化学习研究中一些活跃的前沿方向。

\section*{第 14 章\quad 心理学}

在前面的章节中，我们仅从计算上的考虑出发，发展了各种算法思想。
在本章中，我们从另一个角度来看其中一些算法：心理学的角度，以及心理学关于动物如何学习的研究。
本章的目标有两个：第一，讨论强化学习思想和算法怎样对应于心理学家关于动物学习的发现；第二，解释强化学习正在怎样影响动物学习研究。
强化学习给出的清晰形式化把任务、回报和算法系统化，这正在证明对理解实验数据、提出新类型实验，以及指出哪些因素可能是操纵和测量的关键，都极其有用。
强化学习核心处“在长期上优化回报”的思想，正在帮助我们理解动物学习和行为中一些原本令人困惑的特征。

强化学习与心理学理论之间的某些对应并不令人意外，因为强化学习的发展本来就从心理学学习理论中获得过启发。
不过，如本书所发展出的强化学习，是从人工智能研究者或工程师的视角考察理想化情形，其目标是用高效算法解决计算问题，而不是复制或细致解释动物如何学习。
因此，我们所描述的一些对应关系，把各自领域中独立产生的思想连接了起来。
我们相信，这些接触点特别有意义，因为它们揭示了对学习重要的计算原则，无论学习发生在人工系统还是自然系统中。

在很大程度上，我们描述的是强化学习与一些学习理论之间的对应；这些理论是为了说明大鼠、鸽子和兔子等动物在受控实验室实验中如何学习而发展出来的。
整个 20 世纪进行了成千上万这样的实验，今天仍有许多实验在进行。
虽然这些实验有时被认为与心理学中更广泛的问题无关，但它们探查了动物学习的精细性质，而且常由明确的理论问题推动。
随着心理学把注意力转向行为中更认知的方面，也就是思想和推理等心理过程，动物学习实验在心理学中的作用比过去小了。
但是，这些实验发现了基本而又广泛存在于动物界的学习原则；在设计人工学习系统时，这些原则不应被忽视。
此外，如我们将看到的，认知加工的某些方面也自然地连接到强化学习所提供的计算视角。

本章最后一节包含与我们讨论的联系有关的参考文献，也包含与我们没有讨论的联系有关的参考文献。
我们希望本章能鼓励读者更深入地探究所有这些联系。
最后一节还讨论强化学习中使用的术语如何与心理学术语相联系。
强化学习中的许多术语和短语借自动物学习理论，但这些术语和短语在计算或工程意义上并不总是与它们在心理学中的含义一致。

\subsection*{14.1\quad 预测与控制}

本书描述的算法可分成两大类：用于预测的算法和用于控制的算法。\footnote{这里“控制”的含义不同于动物学习理论中通常的含义；在那里，是环境控制智能体，而不是反过来。参见本章末尾关于术语的说明。}
这两个类别自然地出现在第 3 章给出的强化学习问题求解方法中。
在许多方面，这两个类别分别对应于心理学家广泛研究的两类学习：经典条件化，或巴甫洛夫条件化；以及工具性条件化，或操作性条件化。
由于心理学对强化学习有影响，这些对应并非完全偶然；但它们仍然引人注目，因为它们把来自不同目标的思想连接在一起。

本书中的预测算法估计一些量，这些量取决于智能体环境的特征在未来预期如何展开。
我们特别关注的是估计智能体与环境相互作用时，未来可以预期得到多少奖励。
在这一作用下，预测算法就是策略评估算法，是改进策略的算法中不可缺少的组成部分。
但是，预测算法并不限于预测未来奖励；它们可以预测环境的任何特征。
预测算法与经典条件化之间的对应，建立在共同性质之上：预测即将到来的刺激，无论这些刺激是否具有奖赏性或惩罚性。

工具性条件化，或操作性条件化，实验中的情形则不同。
在这里，实验装置被设置成：动物会根据它做了什么而得到它喜欢的东西（奖励），或得到它不喜欢的东西（惩罚）。
动物学会增加产生受奖赏行为的倾向，并减少产生受惩罚行为的倾向。
强化刺激被说成是依赖于动物行为的；而在经典条件化中，它并不依赖于动物行为（尽管在经典条件化实验中很难去除所有行为依赖性）。
工具性条件化实验类似于启发桑代克效果律的那些实验，我们在第 1 章曾简要讨论过。
控制是这种学习形式的核心，它对应于强化学习中策略改进算法的作用。

把经典条件化理解为预测，把工具性条件化理解为控制，是把我们关于强化学习的计算观点连接到动物学习的出发点；但现实情形比这更复杂。
经典条件化不只是预测；它也涉及动作，因此也是一种控制形式，有时称为巴甫洛夫控制。
此外，经典条件化与工具性条件化会以有趣方式相互作用，在大多数实验情境中，两类学习很可能都参与其中。
尽管有这些复杂性，把经典/工具性区分与预测/控制区分对应起来，仍是连接强化学习和动物学习的一个方便的一阶近似。

在心理学中，术语“强化”用来描述经典条件化和工具性条件化中的学习。
它最初只指一种行为模式的增强，但也常被用于一种行为模式的减弱。
被认为导致行为变化的刺激称为强化物，无论它是否依赖于动物先前的行为。
本章末尾将更详细地讨论这些术语，以及它们与机器学习中术语的关系。

\subsection*{14.2\quad 经典条件化}

著名的俄罗斯生理学家伊万·巴甫洛夫在研究消化系统活动时发现，动物对某些触发刺激的先天反应，会变得可由另一些与先天触发物相当无关的刺激触发。
他的实验对象是狗；这些狗接受过小手术，使其唾液反射的强度能够被准确测量。
在他描述的一个例子中，狗在大多数情形下不会分泌唾液，但在呈现食物后约 5 秒，它会在接下来的几秒内产生大约六滴唾液。
在把另一个与食物无关的刺激，即节拍器的声音，多次安排在食物出现之前之后，狗会像对食物那样，对节拍器声音产生唾液反应。
也就是说，唾液腺的活动被声音冲动调动起来了；而声音本是与食物相当无关的刺激。
巴甫洛夫概括这一发现的重要性时写道：在自然条件下，正常动物显然不仅必须对本身带来即时利益或伤害的刺激作出反应，也必须对其他物理或化学因素作出反应，如声音、光等；这些因素本身只是预示那些刺激的接近。
对较小动物有害的并不是捕食兽的形象和声音本身，而是它的牙齿和爪子。

以这种方式把新刺激连接到先天反射，现在称为经典条件化，或巴甫洛夫条件化。
巴甫洛夫（更准确地说，是他的译者）把先天反应，例如上面示范中的唾液分泌，称为“无条件反应”（UR）；把它们的自然触发刺激，例如食物，称为“无条件刺激”（US）；把由预测性刺激触发的新反应，例如这里同样是唾液分泌，称为“条件反应”（CR）。
一个最初中性的刺激，也就是通常不会诱发强反应的刺激，例如节拍器声音，会随着动物学会它预测 US 而成为“条件刺激”（CS），并因而能对 CS 产生 CR。
这些术语今天仍用于描述经典条件化实验。
US 称为强化物，因为它强化了对 CS 产生 CR 这件事。

两类常见经典条件化实验中刺激的安排如下图所示。
在延迟条件化中，CS 贯穿整个刺激间隔（ISI），也就是 CS 开始与 US 开始之间的时间间隔；图中常见版本还让 CS 在 US 结束时结束。
在迹条件化中，US 在 CS 结束之后才开始，而 CS 结束与 US 开始之间的时间间隔称为迹间隔。

\begin{figure}[htbp]
\centering
\includegraphics[width=0.52\textwidth]{classical-conditioning-timing.png}
\caption*{\textbf{图：}延迟条件化与迹条件化中的刺激时序。}
\end{figure}

巴甫洛夫的狗对节拍器声音分泌唾液，只是经典条件化的一个例子。
经典条件化已经在许多动物物种的许多反应系统中被深入研究。
UR 往往以某种方式具有准备性，例如巴甫洛夫的狗分泌唾液；或以某种方式具有保护性，例如眼睛受到刺激时眨眼，或看到捕食者时冻结不动。
动物在一系列试次中经历 CS--US 预测关系，会学会 CS 预测 US；于是它就能以 CR 对 CS 作出反应，从而为预测到的 US 做准备，或保护自己免受其影响。
有些 CR 与 UR 相似，但开始得更早，并且在提高有效性方面有所不同。
例如，在一种被深入研究的实验中，音调 CS 可靠地预测对兔眼的一股气流（US），气流触发一种 UR，即闭合称为瞬膜的保护性内眼睑。
经过一次或多次试次后，音调会触发一种 CR，即瞬膜闭合；它开始于气流之前，并最终形成合适的时机，使最大闭合恰好发生在气流很可能出现的时候。
由于这种 CR 预期气流而启动并具有适当时机，它比仅仅对刺激性 US 作反应而开始闭合提供更好的保护。
通过学习刺激之间的预测关系而预期重要事件并采取行动，这种能力如此有益，以至于广泛存在于动物界。

\subsection*{14.2.1\quad 阻滞与高阶条件化}

经典条件化实验中观察到了许多有趣性质。
除了 CR 的预期性质之外，两个广泛观察到的性质在经典条件化模型的发展中起了重要作用：阻滞和高阶条件化。
阻滞是指：当一个潜在 CS 与另一个先前已经用来使动物产生该 CR 的 CS 一起呈现时，动物不能学会对这个潜在 CS 产生 CR。
例如，在兔瞬膜条件化的阻滞实验第一阶段，先用音调 CS 和气流 US 训练兔子，使它产生预期气流的瞬膜闭合 CR。
实验第二阶段由额外试次组成，其中第二个刺激，比如光，被加到音调上，形成复合的音调/光 CS，随后仍接同一个气流 US。
在实验第三阶段，仅把第二个刺激，也就是光，呈现给兔子，以检验它是否学会了以 CR 对光作反应。
结果是，兔子对光产生很少甚至没有 CR：对光的学习被先前对音调的学习阻滞了。\footnote{要说明先前对音调的条件化导致了对光学习的阻滞，需要与控制组比较。控制组接受音调/光 CS 的试次，但没有先前对音调的条件化；此时对光的学习不受损。}
这样的阻滞结果挑战了“条件化只依赖简单时间接近”的想法；也就是说，挑战了“US 经常在时间上紧随 CS 出现，就是条件化的必要且充分条件”的想法。
下一节将描述 Rescorla--Wagner 模型，它对阻滞给出了有影响力的解释。

高阶条件化发生在先前已经被条件化的 CS 充当 US，去条件化另一个最初中性的刺激时。
巴甫洛夫描述过一个实验：他的助手首先训练狗对预测食物 US 的节拍器声音分泌唾液。
在这一条件化阶段之后，进行了若干试次：把一块狗最初并不在意的黑色方块放在狗的视线中，随后出现节拍器声音，但之后不再出现食物。
仅仅十个试次后，狗开始仅仅看到黑色方块就分泌唾液，尽管看到方块从未随后接着食物。
节拍器声音本身充当了 US，使黑色方块 CS 获得唾液分泌 CR。
这就是二阶条件化。
如果黑色方块又被用作 US，使另一个原本中性的 CS 建立唾液分泌 CR，那就是三阶条件化，依此类推。
高阶条件化很难展示，尤其是在二阶以上；部分原因是，在高阶条件化试次中，高阶强化物没有被原始 US 反复跟随，因而会失去强化价值。
不过，在适当条件下，例如把一阶试次混入高阶试次，或提供一般性的激活刺激，二阶以上的高阶条件化也可以被展示。
如下文所述，经典条件化的 TD 模型使用了我们方法核心处的自举思想，把 Rescorla--Wagner 模型对阻滞的解释扩展到 CR 的预期性质和高阶条件化。

高阶工具性条件化也会发生。
在这种情况下，一个始终预测初级强化的刺激本身会成为强化物；所谓初级强化，是指其奖赏或惩罚性质由进化内建于动物体内。
预测性刺激会成为次级强化物，或更一般地说，成为高阶强化物或条件强化物；当被预测的强化刺激本身是次级甚至更高阶强化物时，后一术语更好。
条件强化物提供条件强化，即条件奖励或条件惩罚。
条件强化像初级强化一样，会增加动物产生导致条件奖励的行为的倾向，并减少动物产生导致条件惩罚的行为的倾向。

条件强化是一个关键现象。
例如，它解释了为什么我们会为了钱这一条件强化物而工作；钱的价值完全来自“拥有钱”所预测的东西。
在第 13.5 节描述的演员--评论家方法中，评论家使用 TD 方法评估演员的策略，其价值估计向演员提供条件强化，使演员能够改进策略。
这种高阶工具性条件化的类比有助于处理第 1.7 节提到的信用分配问题，因为当初级奖励信号延迟时，评论家仍能逐时刻给演员提供强化。
第 14.4 节将进一步讨论这一点。

\subsection*{14.2.2\quad Rescorla--Wagner 模型}

Rescorla 和 Wagner 创建其模型主要是为了解释阻滞。
Rescorla--Wagner 模型的核心思想是：动物只在事件违反其预期时学习；换言之，只有动物感到“意外”时才学习（这里不一定意味着有意识预期或情绪）。
我们先用他们的术语和记号介绍模型，然后再转到描述 TD 模型时使用的术语和记号。

他们的模型调节复合 CS 中每个成分刺激的“联结强度”；这个数表示该成分以多强或多可靠的方式预测 US。
当由若干成分刺激组成的复合 CS 在经典条件化试次中呈现时，每个成分刺激的联结强度都会改变；这种改变依赖于与整个刺激复合体相关的联结强度，称为“总联结强度”，而不仅依赖各成分本身的联结强度。

设复合 CS 为 $AX$，由成分刺激 $A$ 和 $X$ 组成；动物可能已经经历过刺激 $A$，而刺激 $X$ 可能是新的。
令 $V_A,V_X,V_{AX}$ 分别表示刺激 $A$、$X$ 和复合体 $AX$ 的联结强度。
假设在某试次中，复合 CS $AX$ 后接一个 US，称为刺激 $Y$。
那么刺激成分的联结强度按如下式子改变：
\[
\Delta V_A=\alpha_A\beta_Y(R_Y-V_{AX}),\qquad
\Delta V_X=\alpha_X\beta_Y(R_Y-V_{AX}),
\]
其中 $\alpha_A\beta_Y$ 与 $\alpha_X\beta_Y$ 是步长参数，依赖于 CS 成分和 US 的身份；$R_Y$ 是 US $Y$ 能够支持的联结强度渐近水平。
模型的一个关键假设是，总联结强度 $V_{AX}$ 等于 $V_A+V_X$。
由这些 $\Delta V$ 改变后的联结强度，成为下一试次开始时的联结强度。

为了完整，模型还需要一个反应生成机制，也就是把 $V$ 的值映射到 CR 的方式。
由于这种映射会依赖实验情境的细节，Rescorla 和 Wagner 没有指定具体映射，而只是简单假设：较大的 $V$ 会产生更强或更可能的 CR；负的 $V$ 意味着不会有 CR。

Rescorla--Wagner 模型用一种解释阻滞的方式说明 CR 的获得。
只要刺激复合体的总联结强度 $V_{AX}$ 低于 US $Y$ 能支持的渐近联结强度 $R_Y$，预测误差 $R_Y-V_{AX}$ 就为正。
这意味着，随着试次推进，成分刺激的联结强度 $V_A$ 和 $V_X$ 会增加，直到总联结强度 $V_{AX}$ 等于 $R_Y$；此时联结强度停止变化，除非 US 改变。
当一个新成分被加入到动物已经被条件化过的复合 CS 中时，由于误差已经被降低到零或很小，用增强后的复合体继续条件化几乎不会增加新加入 CS 成分的联结强度。
US 的出现已经几乎被完美预测，因此新 CS 成分几乎不会引入误差或意外。
先前学习阻滞了对新成分的学习。

为了从 Rescorla 和 Wagner 的模型过渡到经典条件化的 TD 模型，我们先用本书一直使用的概念重新表述他们的模型。
具体说，我们匹配第 9.4 节中带线性函数近似的学习记号，并把条件化过程看成：在某试次呈现的复合 CS 基础上，学习预测该试次中“US 的大小”。
US $Y$ 的大小就是上面 Rescorla--Wagner 模型中的 $R_Y$。
我们还引入状态。
因为 Rescorla--Wagner 模型是试次层面的模型，也就是只处理联结强度如何从一个试次变到下一个试次，而不考虑试次内和试次间发生的细节，所以在下一节给出完整 TD 模型之前，我们不必考虑状态在试次内如何变化。
这里我们简单地把状态看成用本试次中存在的 CS 成分集合来标记试次的方式。

因此，假设试次类型或状态 $s$ 由实值特征向量
\[
\mathbf{x}(s)=(x_1(s),x_2(s),\ldots,x_d(s))^\top
\]
描述，其中如果复合 CS 的第 $i$ 个成分 $\mathrm{CS}_i$ 在本试次中存在，则 $x_i(s)=1$，否则为 $0$。
若 $d$ 维联结强度向量为 $\mathbf{w}$，则试次类型 $s$ 的总联结强度为
\begin{equation}
\hat v(s,\mathbf{w})=\mathbf{w}^\top\mathbf{x}(s).
\tag{14.1}
\end{equation}
这对应于强化学习中的价值估计，我们把它看成 US 预测。

现在暂时令 $t$ 表示完整试次编号，而不是通常的时间步编号；当下面扩展到 TD 模型时，我们会恢复 $t$ 的通常含义。
假设 $S_t$ 是对应于试次 $t$ 的状态。
条件化试次 $t$ 把联结强度向量 $\mathbf{w}_t$ 更新为 $\mathbf{w}_{t+1}$：
\begin{equation}
\mathbf{w}_{t+1}=\mathbf{w}_t+\alpha\delta_t\mathbf{x}(S_t),
\tag{14.2}
\end{equation}
其中 $\alpha$ 是步长参数；由于这里描述的是 Rescorla--Wagner 模型，$\delta_t$ 是预测误差
\begin{equation}
\delta_t=R_t-\hat v(S_t,\mathbf{w}_t).
\tag{14.3}
\end{equation}
$R_t$ 是试次 $t$ 上的预测目标，也就是 US 的大小；用 Rescorla 和 Wagner 的术语说，就是该试次的 US 能够支持的联结强度。
注意，由于式（14.2）中有因子 $\mathbf{x}(S_t)$，某试次只会调整该试次中存在的 CS 成分的联结强度。
可以把预测误差看成意外程度的度量，把总联结强度看成动物的预期；当二者与目标 US 大小不匹配时，这个预期被违反。

从机器学习角度看，Rescorla--Wagner 模型是一条误差校正的监督学习规则。
它本质上与最小均方（LMS）或 Widrow--Hoff 学习规则相同：寻找权重，也就是这里的联结强度，使所有误差平方的平均尽可能接近零。
它是一种“曲线拟合”或回归算法，在工程和科学应用中被广泛使用。\footnote{LMS 规则与 Rescorla--Wagner 模型的唯一区别是：LMS 的输入向量成分可以是任意实数，并且至少在最简单版本中，步长参数 $\alpha$ 不依赖输入向量或设定预测目标的刺激身份。}

Rescorla--Wagner 模型在动物学习理论史上非常有影响力，因为它说明一个“机械式”理论可以解释阻滞的主要事实，而不需要诉诸更复杂的认知理论。
例如，它不需要假定动物明确意识到另一个刺激成分被加入，然后向后扫描短时记忆以重新评估涉及 US 的预测关系。
该模型说明，传统的接近理论，也就是认为刺激在时间上接近是学习的必要且充分条件，可以经过简单修正来解释阻滞。
Rescorla--Wagner 模型给出了对阻滞和经典条件化若干其他特征的简单解释，但它不是经典条件化的完整或完美模型。
还有许多不同想法解释其他观察到的效应，对经典条件化诸多精细之处的理解也仍在推进。
下一节描述的 TD 模型同样不是完整或完美的经典条件化模型，但它扩展了 Rescorla--Wagner 模型，处理刺激之间试次内和试次间的时间关系如何影响学习，以及高阶条件化如何产生。

\subsection*{14.2.3\quad TD 模型}

TD 模型是实时模型，而 Rescorla--Wagner 模型是试次层面的模型。
Rescorla--Wagner 模型中的一个步 $t$ 表示整个条件化试次。
该模型不适用于试次进行期间发生的细节，也不适用于试次之间可能发生的事情。
在每个试次内，动物可能经历多种刺激；这些刺激有特定的开始时刻和持续时间。
这些时间关系强烈影响学习。
Rescorla--Wagner 模型也不包含高阶条件化机制；而对 TD 模型来说，高阶条件化是 TD 算法基础处自举思想的自然结果。

为描述 TD 模型，我们从上面对 Rescorla--Wagner 模型的表述开始，但现在 $t$ 标记试次内或试次间的时间步，而不是完整试次。
可以把 $t$ 与 $t+1$ 之间看成一个很小的时间间隔，例如 $0.01$ 秒；把一个试次看成一串状态，每个时间步对应一个状态。
现在步 $t$ 的状态表示刺激在 $t$ 时刻如何被表征的细节，而不仅是某试次中存在的 CS 成分标签。
事实上，我们可以完全放弃试次这个想法。
从动物的角度看，一个试次只是它与世界持续交互经验中的一个片段。
按照我们通常关于智能体与环境交互的观点，可以想象动物正在经历无尽的状态序列 $s$，每个状态都由特征向量 $\mathbf{x}(s)$ 表示。
话虽如此，把试次称为实验中刺激模式重复出现的一段时间，仍常常很方便。

状态特征并不限于描述动物经历的外部刺激；它们也可以描述外部刺激在动物脑中产生的神经活动模式。
这些模式可以依赖历史，也就是由外部刺激序列产生的持久模式。
当然，我们并不知道这些神经活动模式究竟是什么；但 TD 模型这样的实时模型允许人们探索关于外部刺激内部表征的不同假设对学习的影响。
因此，TD 模型不承诺任何特定状态表征。
此外，因为 TD 模型包含跨越刺激之间时间间隔的折扣和资格迹，所以它也能探索折扣和资格迹如何与刺激表征相互作用，从而产生对经典条件化实验结果的预测。

下面将描述一些与 TD 模型一起使用过的状态表征及其含义；但此刻我们先对表征保持中立，只假设每个状态 $s$ 由特征向量
\[
\mathbf{x}(s)=(x_1(s),x_2(s),\ldots,x_n(s))^\top
\]
表示。
那么状态 $s$ 对应的总联结强度仍由式（14.1）给出，与 Rescorla--Wagner 模型相同；但 TD 模型以不同方式更新联结强度向量 $\mathbf{w}$。
当 $t$ 标记时间步而不是完整试次时，TD 模型按如下更新控制学习：
\begin{equation}
\mathbf{w}_{t+1}=\mathbf{w}_t+\alpha\delta_t\mathbf{z}_t,
\tag{14.4}
\end{equation}
它用资格迹向量 $\mathbf{z}_t$ 替代 Rescorla--Wagner 更新（14.2）中的 $\mathbf{x}(S_t)$；并且这里的 $\delta_t$ 不是式（14.3）中的误差，而是 TD 误差：
\begin{equation}
\delta_t=R_{t+1}+\gamma\hat v(S_{t+1},\mathbf{w}_t)-\hat v(S_t,\mathbf{w}_t),
\tag{14.5}
\end{equation}
其中 $\gamma$ 是折扣因子（介于 0 和 1 之间），$R_t$ 是时刻 $t$ 的预测目标，$\hat v(S_{t+1},\mathbf{w}_t)$ 与 $\hat v(S_t,\mathbf{w}_t)$ 是由式（14.1）定义的 $t+1$ 与 $t$ 时刻的总联结强度。

资格迹向量 $\mathbf{z}_t$ 的每个成分 $i$ 会按特征向量 $\mathbf{x}(S_t)$ 的成分 $x_i(S_t)$ 增减，并以由 $\gamma\lambda$ 决定的速率衰减：
\begin{equation}
\mathbf{z}_{t+1}=\gamma\lambda\mathbf{z}_t+\mathbf{x}(S_t).
\tag{14.6}
\end{equation}
这里 $\lambda$ 是通常的资格迹衰减参数。

注意，如果 $\lambda=0$，TD 模型就退化为 Rescorla--Wagner 模型，只有两个例外：其一，$t$ 在两者中的含义不同；其二，在 TD 模型中预测目标 $R$ 有一个时间步的提前。
TD 模型等价于带线性函数近似的半梯度 TD($\lambda$) 算法的后向视图，只是模型中的 $R_t$ 不必像 TD 算法用于学习策略改进的价值函数时那样是奖励信号。

\subsection*{14.2.4\quad TD 模型仿真}

TD 模型这样的实时条件化模型之所以有趣，主要是因为它们能对试次层模型无法表示的广泛情形作出预测。
这些情形涉及可条件化刺激的时序和持续时间、这些刺激相对于 US 的时序，以及 CR 的时序和形状。
例如，US 通常必须在中性刺激开始之后才开始，条件化才会发生；学习速度和有效性还依赖刺激间隔 ISI，也就是 CS 与 US 开始之间的间隔。
当 CR 出现时，它们通常开始于 US 出现之前，并且其时间轮廓在学习过程中发生变化。
在带复合 CS 的条件化中，复合 CS 的成分刺激可能并非同时开始和结束，有时形成所谓的串行复合体，其中成分刺激按时间顺序出现。
这些时序因素使我们必须考虑刺激如何被表征，这些表征如何在试次内和试次间随时间展开，以及它们如何与折扣和资格迹相互作用。

\begin{figure}[htbp]
\centering
\includegraphics[width=0.56\textwidth]{figure-14-1.png}
\caption*{\textbf{图 14.1：}TD 模型有时使用的三种刺激表征（按列展示）：完整串行复合体、微刺激和存在表征。每一行表示刺激表征中的一个元素。三种表征在时间泛化梯度上不同：完整串行复合体在邻近时间点之间没有泛化；存在表征在邻近时间点之间完全泛化；微刺激表征介于两者之间。时间泛化程度决定了学习 US 预测时的时间粒度。}
\end{figure}

图 14.1 展示了探索 TD 模型行为时用过的三种刺激表征：完整串行复合体（CSC）、微刺激（MS）和存在表征。
这些表征的差异在于，当刺激存在时，它们在多大程度上强制邻近时间点之间发生泛化。

图 14.1 中最简单的表征是右列的存在表征。
这种表征对试次中存在的每个 CS 成分只有一个特征；只要该成分存在，该特征取值为 1，否则为 0。\footnote{在我们的形式化中，一个试次中的每个时间步 $t$ 都有不同状态 $S_t$。若复合 CS 包含 $n$ 个在试次中不同时间、以不同持续时间出现的成分 CS，则每个成分 $\mathrm{CS}_i$ 有一个特征 $x_i$；当 $\mathrm{CS}_i$ 存在时 $x_i(S_t)=1$，否则为 0。}
存在表征并不是关于动物脑中刺激如何被表征的现实假设；但如下文所述，TD 模型配合这种表征能产生许多经典条件化中看到的时序现象。

对 CSC 表征（图 14.1 左列）来说，每个外部刺激的开始都会引发一串精确定时、持续时间很短的内部信号，这串信号持续到外部刺激结束。\footnote{在我们的形式化中，对试次中存在的每个 CS 成分 $\mathrm{CS}_i$，以及试次中的每个时间步 $t$，都有单独特征 $x_i^t$；若 $t=t'$ 且 $\mathrm{CS}_i$ 在 $t'$ 时刻存在，则 $x_i^t(S_{t'})=1$，否则为 0。}
这就像假设动物神经系统有一只钟，能在刺激呈现期间精确追踪时间；工程师称之为“抽头延迟线”。
与存在表征一样，CSC 表征作为关于脑如何内部表征刺激的假设并不现实；但它是一种“有用的虚构”，因为它能揭示 TD 模型在较少受刺激表征约束时如何工作。
CSC 表征也用于大多数脑内多巴胺产生神经元的 TD 模型；这是第 15 章将讨论的主题。
人们常把 CSC 表征看作 TD 模型的必要组成部分，尽管这种看法是错误的。

MS 表征（图 14.1 中列）与 CSC 表征相似：每个外部刺激都会引发一串内部刺激。
但在这种情况下，内部刺激，也就是微刺激，并不是如此有限且互不重叠；它们在时间上延展并重叠。
随着刺激开始后的时间流逝，不同微刺激集合变得或多或少活跃，每个后续微刺激在时间上逐渐变宽，并达到更低的最大水平。
当然，取决于微刺激的性质，可以有许多 MS 表征；文献中已经研究了若干例子，其中有些还提出动物脑可能怎样生成它们。
作为关于刺激神经表征的假设，MS 表征比存在表征或 CSC 表征更现实，而且允许 TD 模型的行为与动物实验中观察到的更广泛现象联系起来。
特别是，如果假设微刺激级联不仅由 CS 引发，也由 US 引发，并研究微刺激、资格迹和折扣之间相互作用对学习的显著影响，那么 TD 模型有助于形成假说，解释许多经典条件化的精细现象以及动物脑可能如何产生它们。
关于这一点，下面还会说更多，尤其是在第 15 章讨论强化学习和神经科学时。

即使使用简单的存在表征，TD 模型也能产生 Rescorla--Wagner 模型所解释的经典条件化基本性质，并且还能产生超出试次层模型范围的条件化特征。
例如，如前所述，经典条件化的一个显著特征是：US 通常必须在中性刺激开始之后才开始，条件化才会发生；并且在条件化之后，CR 会在 US 出现之前开始。
换言之，条件化通常需要正的 ISI，而 CR 通常预期 US。
条件化强度（例如由 CS 诱发的 CR 百分比）如何依赖 ISI，会随物种和反应系统而有很大差异，但通常具有以下性质：当 ISI 为零或负时，也就是 US 开始与 CS 开始同时或更早时，条件化可忽略；它在某个正 ISI 上增加到最大，此时条件化最有效；之后又在某个随反应系统差异很大的间隔后下降到零。
TD 模型中这种依赖的精确形状取决于参数值和刺激表征细节，但 ISI 依赖性的这些基本特征是 TD 模型的核心性质。

串行复合条件化，也就是复合 CS 的成分按顺序出现时，产生的一个理论问题，是远程联结的促进。
研究发现，如果第一个 CS（CSA）与 US 之间的空迹间隔被第二个 CS（CSB）填充，从而形成串行复合刺激，那么对 CSA 的条件化会被促进。
下图显示了使用存在表征的 TD 模型在这类实验仿真中的行为；其上方显示了时序细节。
与实验结果一致，模型显示：由于第二个 CS 的存在，第一个 CS 的条件化速度和渐近条件化水平都得到促进。

\begin{figure}[htbp]
\centering
\includegraphics[width=0.58\textwidth]{remote-associations-td.png}
\caption*{\textbf{图：}TD 模型中远程联结的促进。}
\end{figure}

Egger 和 Miller 的一个著名实验展示了试次内刺激之间的时间关系对条件化的影响。
该实验包含两个以延迟构型重叠的 CS，如下图上方所示。
尽管 CSB 与 US 的时间关系更好，CSA 的存在却显著降低了对 CSB 的条件化，相比之下，在 CSA 缺席的控制条件中则不会如此。
下图直接显示了 TD 模型在使用存在表征仿真该实验时生成的同样结果。

\begin{figure}[htbp]
\centering
\includegraphics[width=0.58\textwidth]{egger-miller-td.png}
\caption*{\textbf{图：}TD 模型中的 Egger--Miller 效应。}
\end{figure}

TD 模型解释阻滞，是因为它像 Rescorla--Wagner 模型一样，是一条误差校正学习规则。
不过，除了能解释基本阻滞结果之外，TD 模型还预测：如果被阻滞的刺激在时间上前移，使其开始早于阻滞刺激的开始，那么阻滞会被反转。
这个行为特征值得注意，因为模型提出时尚未观察到它。
回忆阻滞情形：如果动物已经学会一个 CS 预测 US，那么对新加入的第二个 CS 也预测 US 的学习会大幅减少，也就是被阻滞。
但是，如果新加入的第二个 CS 比预训练 CS 更早开始，那么按 TD 模型，对新加入 CS 的学习不会被阻滞。
事实上，随着训练继续，新加入 CS 获得联结强度，预训练 CS 会失去联结强度。
TD 模型在这些条件下的行为见图 14.2。
该仿真实验与前页的 Egger--Miller 实验不同：其中较短且开始较晚的 CS 先接受预训练，直到它与 US 完全联结。
这一出人意料的预测促使 Kehoe、Schreurs 和 Graham 使用已被深入研究的兔瞬膜制备进行实验。
结果证实了模型预测；他们还指出，非 TD 模型很难解释这些数据。

\begin{figure}[htbp]
\centering
\includegraphics[width=0.58\textwidth]{figure-14-2.png}
\caption*{\textbf{图 14.2：}TD 模型中，时间优先性压倒阻滞。}
\end{figure}

在 TD 模型中，较早的预测刺激优先于较晚的预测刺激，因为与本书描述的所有预测方法一样，TD 模型基于备份或自举思想：对联结强度的更新会把某个状态处的强度推向较后状态处的强度。
自举的另一个结果是，TD 模型能解释高阶条件化；这是 Rescorla--Wagner 模型和类似模型范围之外的经典条件化特征。
如上所述，高阶条件化是指先前被条件化的 CS 可以充当 US，来条件化另一个最初中性的刺激。
图 14.3 展示了 TD 模型（仍使用存在表征）在高阶条件化实验中的行为；这里是二阶条件化。
在第一阶段（图中未示出），CSB 被训练来预测 US，因此它的联结强度增加，这里达到 1.65。
在第二阶段，在没有 US 的情况下，CSA 与 CSB 按图顶部所示的顺序安排配对。
尽管 CSA 从未与 US 配对，它仍获得联结强度。
随着训练继续，CSA 的联结强度达到峰值后又下降；这是因为作为次级强化物的 CSB 的联结强度下降，因而失去提供次级强化的能力。
CSB 的联结强度下降，是因为 US 在这些高阶条件化试次中没有出现。
对 CSB 来说，这些是消退试次，因为它与 US 的预测关系被破坏，使其作为强化物的能力下降。
动物实验中也能看到同样模式。
高阶条件化试次中条件强化的这种消退，使得如果不偶尔插入一阶试次、周期性刷新原始预测关系，就很难展示高阶条件化。

\begin{figure}[htbp]
\centering
\includegraphics[width=0.58\textwidth]{figure-14-3.png}
\caption*{\textbf{图 14.3：}TD 模型中的二阶条件化。}
\end{figure}

TD 模型之所以产生二阶和高阶条件化的类比，是因为 TD 误差（14.5）中出现了 $\gamma\hat v(S_{t+1},\mathbf{w}_t)-\hat v(S_t,\mathbf{w}_t)$。
这意味着，由于先前学习，$\gamma\hat v(S_{t+1},\mathbf{w}_t)$ 可以不同于 $\hat v(S_t,\mathbf{w}_t)$，从而使 $\delta_t$ 非零，也就是产生时间差分。
在式（14.5）中，这个差分与 $R_{t+1}$ 具有同样地位；这意味着就学习而言，一个时间差分与一个 US 的发生没有区别。
事实上，TD 算法这一特征是它得以发展的主要原因之一；我们现在通过第 6 章所述它与动态规划的联系来理解这一点。
价值的自举与二阶及更高阶条件化密切相关。

在上面描述的 TD 模型行为例子中，我们只考察了 CS 成分联结强度的变化。
我们还没有考察模型如何预测动物 CR 的性质：它们的时序、形状，以及它们如何在条件化试次中发展。
这些性质依赖物种、被观察的反应系统以及条件化试次的参数；但在许多不同动物和不同反应系统的实验中，CR 的幅度或 CR 的概率会随着 US 的预期时间临近而增加。
例如，在前面提到的兔瞬膜反应经典条件化中，随着条件化试次推进，从 CS 开始到瞬膜开始移过眼睛之间的延迟逐渐缩短；这种预期性闭合的幅度在 CS 与 US 之间的间隔内逐渐增加，直到瞬膜在 US 的预期时间达到最大闭合。
这种 CR 的时序和形状对其适应意义至关重要：过早遮住眼睛会降低视觉（即使瞬膜是半透明的），过晚遮住眼睛则几乎没有保护价值。
捕捉这些 CR 特征，对经典条件化模型来说是有挑战的。

TD 模型的定义本身不包含任何机制，用来把 US 预测的时间过程 $\hat v(S_t,\mathbf{w}_t)$ 转换成可与动物 CR 性质比较的轮廓。
最简单选择是让仿真 CR 的时间过程等于 US 预测的时间过程。
在这种情况下，仿真 CR 的特征及其随试次的变化，只依赖于所选刺激表征以及模型参数 $\alpha,\gamma,\lambda$ 的取值。

图 14.4 展示了使用图 14.1 中三种表征时，学习过程中不同点处 US 预测的时间过程。
这些仿真中，US 在 CS 开始后 25 个时间步发生，$\alpha=.05$，$\lambda=.95$，$\gamma=.97$。
使用 CSC 表征时（图 14.4 左），TD 模型形成的 US 预测曲线在 CS 与 US 之间的间隔中指数式上升，直到 US 发生时恰好达到最大值（时间步 25）。
这种指数式上升是 TD 模型学习规则中折扣的结果。
使用存在表征时（图 14.4 中），刺激存在期间 US 预测几乎恒定，因为每个刺激只有一个权重或联结强度需要学习。
因此，带存在表征的 TD 模型不能重现 CR 时序的许多特征。
使用 MS 表征时（图 14.4 右），TD 模型的 US 预测发展更复杂。
经过 200 个试次后，预测轮廓已经合理近似于 CSC 表征产生的 US 预测曲线。

\begin{figure}[htbp]
\centering
\includegraphics[width=0.92\textwidth]{figure-14-4.png}
\caption*{\textbf{图 14.4：}TD 模型在获得过程中、使用三种不同刺激表征时的 US 预测时间过程。左：完整串行复合体使 US 预测在间隔内指数式增加，并在 US 时间达到峰值。中：存在表征使 US 预测收敛到近似常数水平。右：微刺激表征在渐近时通过不同微刺激的线性组合，近似 CSC 中指数式上升的时间过程。}
\end{figure}

图 14.4 中的 US 预测曲线并不旨在精确匹配任何特定动物实验中 CR 发展出的轮廓；它们说明刺激表征对 TD 模型导出的预测有强烈影响。
此外，虽然这里只能提到，刺激表征如何与折扣和资格迹相互作用，对决定 TD 模型产生的 US 预测轮廓性质很重要。
另一个超出此处讨论范围的维度，是不同反应生成机制的影响；这些机制把 US 预测转换成 CR 轮廓。
图 14.4 中的轮廓是“原始”的 US 预测轮廓。
不过，即使不作关于动物脑如何从 US 预测产生外显反应的特殊假设，图 14.4 中 CSC 和 MS 表征的轮廓也会随着 US 时间临近而增加，并在 US 时间达到最大值，这与许多动物条件化实验所见一致。

TD 模型在与特定刺激表征和反应生成机制结合时，能够解释动物经典条件化实验中观察到的惊人广泛的现象；但它远非完美模型。
要生成经典条件化的其他细节，模型还需要扩展，也许要加入基于模型的要素，以及自适应改变某些参数的机制。
另一些经典条件化建模方法则明显不同于 Rescorla--Wagner 式误差校正过程。
例如，贝叶斯模型在概率框架内工作，在其中经验修正概率估计。
所有这些模型都对理解经典条件化有用。

TD 模型也许最值得注意的特征是，它基于一种理论，也就是本书描述的理论。
这为动物神经系统在条件化过程中试图做什么提供了一种解释：它试图形成准确的长期预测，但要受刺激如何被表征和神经系统如何工作的限制。
换言之，它提示了经典条件化的一种规范性解释，其中长期预测而不是即时预测是关键特征。

经典条件化 TD 模型的发展，是一个明确目标为建模动物学习行为若干细节的例子。
因此，TD 学习除了作为一种算法，也成为这个生物学习某些方面模型的基础。
如第 15 章将讨论的，TD 学习还被证明是一个有影响力模型的基础，该模型解释产生多巴胺的神经元活动；多巴胺是哺乳动物脑中的一种化学物质，深度参与奖励加工。
这些都是强化学习理论与动物行为和神经数据发生细致接触的例子。

现在我们转向强化学习与工具性条件化实验中的动物行为之间的对应；工具性条件化是动物学习心理学家研究的另一大类实验。

\subsection*{14.3\quad 工具性条件化}

在工具性条件化实验中，学习依赖行为的结果：强化刺激的给出依赖动物做了什么。
相反，在经典条件化实验中，强化刺激即 US 的给出独立于动物行为。
工具性条件化通常被认为与操作性条件化相同；后者是 B. F. Skinner 为行为依赖强化的实验引入的术语。
不过，使用这两个术语的研究者在实验和理论上有若干差异，下面会触及其中一些。
我们将专门使用“工具性条件化”一词，指强化依赖于行为的实验。
工具性条件化的根源可追溯到美国心理学家 Edward Thorndike 在本书第一版出版前一百年左右进行的实验。

Thorndike 观察猫被放入“迷箱”时的行为；猫可以通过适当动作从这样的箱子中逃出。
例如，一只猫可以通过执行三个分离动作的序列打开某个箱子的门：按下箱子后部的平台，用爪拉绳，以及把一根杆向上或向下推。
当猫第一次被放入迷箱、而食物在外面可见时，除了少数之外，Thorndike 的猫都表现出“明显的不适迹象”和异常剧烈的活动，“本能地挣扎着逃离禁闭”。

\begin{figure}[htbp]
\centering
\includegraphics[width=0.52\textwidth]{thorndike-puzzle-box.png}
\caption*{\textbf{图：}Thorndike 的一个迷箱。}
\end{figure}

在使用不同猫和具有不同逃脱机制的箱子的实验中，Thorndike 记录每只猫在每个箱子中多次经历后逃出所需的时间。
他观察到，所需时间几乎总是随连续经验而减少，例如从 300 秒减少到 6 或 7 秒。
他这样描述迷箱中猫的行为：在冲动挣扎中到处乱抓箱子的猫，很可能会抓到绳、环或按钮而打开门。
逐渐地，所有其他不成功的冲动会被淘汰，而通向成功动作的特定冲动会被由此产生的愉悦巩固下来；直到许多试次之后，猫再被放入箱子时，会立即以确定方式抓按钮或环。

这些以及其他实验（其中一些涉及狗、小鸡、猴子甚至鱼）促使 Thorndike 提出若干学习“定律”，最有影响力的是效果律；我们在第 1 章曾引用过它的一个版本。
这一定律描述了通常所谓的试错学习。
如第 1 章所说，效果律的许多方面曾引发争议，其细节也多年来不断修改。
然而，这一定律以某种形式表达了一个持久的学习原则。

强化学习算法的基本特征对应于效果律所描述的动物学习特征。
第一，强化学习算法是选择性的；也就是说，它们尝试各种备选项，并通过比较后果在其中选择。
第二，强化学习算法是联结性的；也就是说，由选择找到的备选项与特定情境或状态相联结，形成智能体的策略。
像效果律描述的学习一样，强化学习不只是寻找产生大量奖励的动作，也是在把这些动作连接到情境或状态。
Thorndike 使用“选择并连接”来描述学习。
进化中的自然选择是选择过程的典型例子，但它并不是联结性的（至少按通常理解如此）；监督学习是联结性的，但不是选择性的，因为它依赖直接告诉智能体如何改变行为的指令。

用计算术语说，效果律描述了一种把搜索和记忆结合起来的基本方式：搜索表现为在每种情境中尝试许多动作并从中选择；记忆表现为把情境与到目前为止在这些情境中效果最好的动作联结起来。
搜索和记忆是所有强化学习算法的必要组成部分，无论记忆采取智能体策略、价值函数还是环境模型的形式。

强化学习算法需要搜索，意味着它必须以某种方式探索。
动物显然也会探索；早期动物学习研究者曾争论，在像 Thorndike 迷箱这样的情境中，动物选择动作时受到多大程度的引导。
动作是“绝对随机的、盲目的摸索”的结果，还是存在某种程度的引导，无论来自先前学习、推理或其他手段？
包括 Thorndike 在内的一些思想家似乎采取前一种立场，另一些人则更偏向有意探索。
强化学习算法允许智能体在选择动作时使用不同程度的引导。
本书算法中使用的探索形式，例如 $\varepsilon$-贪心和上置信界动作选择，只是最简单的几种。
更复杂的方法是可能的；唯一要求是，算法要有效工作，必须存在某种探索形式。

我们的强化学习处理方式允许任意时刻可用动作集合依赖于环境当前状态；这呼应了 Thorndike 在猫的迷箱行为中观察到的东西。
猫从它们在当前情境中本能会执行的动作中选择动作，Thorndike 称这些为“本能冲动”。
第一次被放入迷箱时，猫会极有力地抓、挠和咬；这些是猫发现自己处于受限空间时的本能反应。
成功动作从这些动作中被选择出来，而不是从所有可能动作或活动中选择出来。
这类似于我们形式化中的特征：从状态 $s$ 选择的动作属于可容许动作集合 $\mathcal{A}(s)$。
指定这些集合是强化学习的重要方面，因为它能极大简化学习。
它们类似于动物的本能冲动。
另一方面，Thorndike 的猫也可能不是只从本能冲动集合中选择，而是按某种本能的、依赖情境的动作顺序进行探索。
这也是使强化学习更容易的一种方式。

受效果律影响最显著的动物学习研究者包括 Clark Hull 和 B. F. Skinner。
他们研究的中心思想是根据行为后果选择行为。
强化学习与 Hull 的理论有一些共同特征；后者包含类似资格的机制和次级强化，用以解释当动作和随后强化刺激之间存在显著时间间隔时仍能学习的能力。
随机性也通过 Hull 所谓“行为振荡”在其理论中发挥作用，用以引入探索行为。

Skinner 并不完全接受效果律中的记忆方面。
由于反感联结链路的思想，他转而强调从自发发出的行为中选择。
他引入“操作性”一词，以强调动作对动物环境产生的效果的关键作用。
不同于 Thorndike 等人的实验由一串分离试次组成，Skinner 的操作性条件化实验允许动物被试长时间不间断地行为。
他发明了操作性条件化箱，现在称为“Skinner 箱”；最基本版本包含一根杠杆或一个键，动物可按压以获得食物或水等奖励，而奖励按照明确规则给出，这个规则称为强化时程。
通过把杠杆按压的累计次数记录为时间函数，Skinner 及其追随者能够研究不同强化时程对动物按压杠杆速率的影响。
用本书给出的强化学习原则来建模这类实验结果尚不成熟，但本章末尾的文献和历史备注会提到一些例外。

Skinner 的另一项贡献来自他认识到一种训练动物的有效方法：强化目标行为的逐次近似，他称这一过程为塑造。
虽然包括 Skinner 本人在内的其他人也用过这种技术，但当他和同事试图训练鸽子用喙击打一只木球来打保龄球时，这种技术的重要性给他留下了深刻印象。
他们等了很久也看不到任何可以强化的击打动作，于是决定强化任何与击打稍有相似的反应，也许一开始只是看着球；然后再选择越来越接近最终形式的反应。
结果使他们惊讶：几分钟内，球就在箱壁之间反弹，就像鸽子是壁球冠军一样。
鸽子不仅学会了对鸽子来说不寻常的行为，而且是在一种交互过程中快速学会的；在这个过程中，它的行为和强化依赖关系相互响应地变化。
Skinner 把改变强化依赖关系的过程比作雕塑家把黏土塑造成目标形状。
塑造对计算强化学习系统同样是强有力的技术。
当智能体很难接收到任何非零奖励信号时，不论因为有奖励的情境稀疏，还是因初始行为无法到达这些情境，从一个更容易的问题开始，并随着智能体学习逐步增加难度，可能是一种有效甚至不可或缺的策略。

在工具性条件化语境中特别相关的心理学概念是动机，它指影响行为方向和强度（或活力）的过程。
例如，Thorndike 的猫有动机逃出迷箱，因为它们想得到箱外的食物。
获得这个目标对它们有奖赏性，并强化了让它们逃脱的动作。
动机概念有许多维度，很难精确连接到强化学习的计算视角，但其中一些维度与强化学习有清楚联系。

在一种意义上，强化学习智能体的奖励信号位于其动机基础处：智能体有动机最大化长期接收的总奖励。
因此，动机的一个关键方面是，什么使智能体的经验具有奖赏性。
在强化学习中，奖励信号依赖于强化学习智能体环境的状态和智能体动作。
此外，如第 1 章指出的，智能体环境状态不仅包括承载智能体的机器、机体或机器人外部的信息，也包括该机器内部的信息。
某些内部状态成分对应于心理学家所说的动物动机状态，这会影响什么对动物有奖赏性。
例如，饥饿动物从进食中获得的奖赏，比刚吃过一顿满足餐的动物更大。
状态依赖这一概念足够宽泛，能容纳许多类型的调制性影响来生成奖励信号。

价值函数给出了与心理学家动机概念的进一步联系。
如果选择动作的最基本动机是尽可能多地获得奖励，那么对使用价值函数选择动作的强化学习智能体来说，更近端的动机就是沿其价值函数梯度上升，也就是选择预期通向最高价值下一状态的动作（或本质上等价地，选择具有最大动作价值的动作）。
对这些智能体来说，价值函数是决定其行为方向的主要驱动力。

动机的另一个维度是，动物的动机状态不仅影响学习，也影响学习之后行为的强度或活力。
例如，在学会在迷宫目标箱中找到食物之后，饥饿大鼠会比不饥饿大鼠更快跑向目标箱。
动机的这一方面与本书给出的强化学习框架连接得不那么直接，但本章末尾会引用若干基于强化学习的行为活力理论。

现在转向当强化刺激发生在被强化事件之后相当久时的学习问题。
强化学习算法用来支持延迟强化下学习的机制，即资格迹和 TD 学习，与心理学家关于动物如何在这些条件下学习的假设密切对应。

\subsection*{14.4\quad 延迟强化}

效果律需要对联结产生向后的影响；一些早期批评者无法想象现在如何影响过去的东西。
这种担忧因一个事实而加剧：即使动作与随后的奖励或惩罚之间存在相当大的延迟，学习仍能发生。
类似地，在经典条件化中，当 US 开始在 CS 结束后经过不可忽略的时间间隔才出现时，学习也能发生。
我们把这称为延迟强化问题，它与 Minsky 所谓“学习系统的信用分配问题”相关：如何把成功的功劳分配给可能参与产生成功的许多决策？
本书介绍的强化学习算法包含两个处理这一问题的基本机制。
第一是使用资格迹；第二是使用 TD 方法学习价值函数，使其对动作提供几乎即时的评价（在工具性条件化实验这样的任务中），或提供即时预测目标（在经典条件化实验这样的任务中）。
这两种方法都对应于动物学习理论中提出的类似机制。

巴甫洛夫指出，每个刺激都必定在神经系统中留下某种迹，并在刺激结束后持续一段时间。
他提出，当 CS 结束与 US 开始之间存在时间间隙时，刺激迹使学习成为可能。
只要 US 到来时仍有 CS 的迹，学习就通过迹与 US 的同时存在而发生。
在这些条件下的条件化至今仍称为迹条件化。
第 15 章会讨论神经系统中若干可能的迹机制。

刺激迹也被提出用于桥接工具性条件化中动作与随后的奖励或惩罚之间的时间间隔。
例如，在 Hull 有影响力的学习理论中，“整体刺激迹”解释了他所谓动物的目标梯度；目标梯度描述工具性条件化反应的最大强度如何随强化延迟增加而降低。
Hull 假设动物动作留下内部刺激，其迹按动作发生后的时间指数衰减。
根据当时可用的动物学习数据，他假设这些迹在 30 到 40 秒后实际上降为零。

本书算法中使用的资格迹类似于 Hull 的迹：它们是过去状态访问或过去状态--动作对的衰减迹。
Klopf 在其神经元理论中引入资格迹，把它们看作突触过去活动在时间上延展的迹；突触是神经元之间的连接。
Klopf 的迹比我们算法使用的指数衰减迹更复杂；第 15.9 节讨论其理论时会进一步说明。

为解释超过刺激迹跨度的更长时间目标梯度，Hull 提出较长梯度来自从目标向后传递的条件强化，这一过程与他的整体刺激迹共同作用。
动物实验表明，如果条件有利于延迟期间条件强化的发展，学习随延迟增加而下降的程度，会小于那些阻碍次级强化的条件。
如果延迟间隔中有规律出现的刺激，条件强化就会受到促进。
于是就好像奖励并未真正延迟，因为存在更即时的条件强化。
因此，Hull 想象存在一个基于由刺激迹中介的初级强化延迟的初级梯度，并且它被条件强化逐渐修改和延长。

本书中同时使用资格迹和价值函数来支持延迟强化下学习的算法，对应于 Hull 关于动物如何在这些条件下学习的假说。
第 13.5、15.7 和 15.8 节讨论的演员--评论家架构最清楚地说明了这种对应。
评论家使用 TD 算法学习与系统当前行为相关的价值函数，也就是预测当前策略的回报。
演员基于评论家的预测，更准确地说，基于评论家预测的变化，来更新当前策略。
评论家产生的 TD 误差作为演员的条件强化信号，即使初级奖励信号本身显著延迟，也能提供即时的表现评价。
估计动作价值函数的算法，例如 Q 学习和 Sarsa，也类似地使用 TD 学习原则，通过条件强化支持延迟强化下的学习。
第 15 章讨论的 TD 学习与多巴胺产生神经元活动之间的密切平行，为强化学习算法与 Hull 学习理论这一方面之间的联系提供了额外支持。

\subsection*{14.5\quad 认知地图}

基于模型的强化学习算法使用环境模型；这些模型与心理学家所谓认知地图有共同要素。
回忆第 8 章关于规划与学习的讨论：环境模型指智能体可用来预测环境如何响应其动作的任何东西，预测内容包括状态转移和奖励；规划指从这样的模型计算策略的任何过程。
环境模型由两部分组成：状态转移部分编码动作对状态变化影响的知识；奖励模型部分编码每个状态或每个状态--动作对预期奖励信号的知识。
基于模型的算法通过使用模型来预测可能行动路线的后果，选择动作；这些后果以未来状态和预期从这些状态产生的奖励信号来表示。
最简单的规划就是比较若干“想象的”决策序列集合的预测后果。

动物是否使用环境模型；如果使用，这些模型是什么样、如何被学习；这些问题在动物学习研究史上曾发挥重要作用。
一些研究者通过展示潜伏学习，挑战了当时占主导的刺激--反应（S--R）学习和行为观点；这种观点对应于学习策略的最简单无模型方式。
在最早的潜伏学习实验中，两组大鼠在迷宫中奔跑。
对实验组来说，实验第一阶段没有奖励；第二阶段开始时，食物突然被放入迷宫目标箱。
对控制组来说，两个阶段中目标箱一直有食物。
问题是，实验组大鼠在第一阶段没有食物奖励时是否学到了什么。
虽然实验大鼠在第一阶段这个无奖励阶段似乎没有学到太多，但一旦它们发现第二阶段引入的食物，就迅速赶上控制组大鼠。
结论是，在无奖励时期，实验组大鼠正在形成关于迷宫的潜伏学习；一旦引入奖励，它们就能利用这种学习。

潜伏学习与心理学家 Edward Tolman 关系最密切。
Tolman 把这一结果以及类似结果解释为：动物能够在没有奖励或惩罚的情况下学习“环境的认知地图”，并在后来有动机达到目标时使用该地图。
认知地图也允许大鼠规划一条不同于最初探索时走过路线的通往目标的路线。
这类结果的解释导致了行为主义/认知二分法核心处持久的争议。
用现代术语说，认知地图并不限于空间布局模型，而更一般地是环境模型，或动物“任务空间”的模型。
用认知地图解释潜伏学习实验，类似于声称动物使用基于模型的算法，并且环境模型即使没有明确奖励或惩罚也能被学习。
当动物被奖励或惩罚的出现所激励时，模型随后用于规划。

Tolman 对动物如何学习认知地图的解释是：动物在探索环境时通过经历刺激序列，学习刺激--刺激（S--S）联结。
在心理学中，这称为期待理论：给定 S--S 联结，一个刺激的发生会产生关于下一个刺激的期待。
这很像控制工程师所谓系统辨识，即从带标签训练样例中学习一个未知动力系统的模型。
在最简单的离散时间版本中，训练样例是 $S$--$S'$ 对，其中 $S$ 是状态，$S'$ 是后续状态标签。
观察到 $S$ 时，模型产生将会观察到 $S'$ 的“期待”。
对规划更有用的模型还涉及动作，因此样例形如 $SA$--$S'$，即在状态 $S$ 执行动作 $A$ 时预期 $S'$。
学习环境如何产生奖励也很有用。
此时样例形如 $S$--$R$ 或 $SA$--$R$，其中 $R$ 是与 $S$ 或 $SA$ 对相关的奖励信号。
这些都是监督学习形式，通过它们，智能体无论在探索环境时是否接收到非零奖励信号，都能获得类似认知地图的模型。

\subsection*{14.6\quad 习惯性行为与目标导向行为}

无模型和基于模型的强化学习算法之间的区别，对应于心理学家在习得行为模式的习惯性控制和目标导向控制之间所作的区别。
习惯是由适当刺激触发、随后或多或少自动执行的行为模式。
按心理学家的用法，目标导向行为是有目的的；它受关于目标价值以及动作与其后果关系的知识控制。
习惯有时被说成受先行刺激控制，而目标导向行为则被说成受其后果控制。
目标导向控制的优点是，当环境改变对动物动作的反应方式时，它能迅速改变动物行为。
习惯性行为能快速响应熟悉环境的输入，但无法快速适应环境变化。
目标导向行为控制的发展，很可能是动物智能进化中的一次重大进步。

图 14.5 说明了无模型和基于模型的决策策略在一个假想任务中的差异。
在该任务中，一只大鼠必须穿过一个迷宫；迷宫具有可区分的目标箱，每个目标箱提供图中所示大小的奖励。
从 $S_1$ 开始，大鼠必须先选择左（L）或右（R），然后在 $S_2$ 或 $S_3$ 再选择 L 或 R，以到达某个目标箱。
目标箱是大鼠分幕任务每一幕的终止状态。
无模型策略依赖已存储的状态--动作对价值。
这些动作价值估计了大鼠从每个非终止状态采取每个动作可预期获得的最高回报。
它们通过多次从起点到终点跑迷宫获得。
当动作价值已经足够好地估计最优回报时，大鼠只需在每个状态选择动作价值最大的动作，就能作出最优决策。

\begin{figure}[htbp]
\centering
\includegraphics[width=0.82\textwidth]{figure-14-5.png}
\caption*{\textbf{图 14.5：}解决一个假想序贯动作选择问题的基于模型策略与无模型策略。上：大鼠在具有不同目标箱的迷宫中导航，每个目标箱关联图示奖励。左下：无模型策略依赖许多学习试次中获得的状态--动作对存储动作价值。右下：基于模型策略学习环境模型，包括状态--动作--下一状态转移知识和目标箱奖励模型；大鼠可用模型模拟动作选择序列，寻找回报最高的路径。}
\end{figure}

在图 14.5 所示情形中，当动作价值估计足够准确时，大鼠从 $S_1$ 选择 L，从 $S_2$ 选择 R，得到最大回报 4。
另一种无模型策略可能直接依赖缓存策略而不是动作价值，直接建立从 $S_1$ 到 L、从 $S_2$ 到 R 的联系。
这两种策略都不依赖环境模型。
既不需要查阅状态转移模型，也不需要在目标箱特征与其提供的奖励之间建立连接。

图 14.5 右下说明基于模型的策略。
它使用由状态转移模型和奖励模型组成的环境模型。
状态转移模型画成决策树；奖励模型把目标箱的可区分特征与其中可获得的奖励相联系。
（与 $S_1,S_2,S_3$ 相关的奖励也是奖励模型的一部分，但这里它们为零，故未画出。）
基于模型的智能体可以用模型模拟动作选择序列，寻找产生最高回报的路径，从而决定在每个状态向哪边转。
在本例中，回报就是路径末端结果得到的奖励。
若模型足够准确，大鼠会先选择 L 再选择 R，以得到奖励 4。
比较被模拟路径的预测回报是一种简单规划形式；第 8 章讨论过规划可以有多种方式。

当无模型智能体的环境改变对其动作的反应方式时，智能体必须在改变后的环境中获得新经验，并在此期间更新其策略和/或价值函数。
例如，在图 14.5 左下的无模型策略中，如果某个目标箱不知为何改为提供不同奖励，大鼠必须穿过迷宫，可能要多次，才能在到达该目标箱时经历新奖励，同时基于这种经验更新策略或动作价值函数（或两者）。
关键点是，对无模型智能体来说，要改变策略在某状态指定的动作，或改变与某状态相关的动作价值，它必须移动到该状态，从那里行动，可能还要多次，并经历动作的后果。

基于模型的智能体则能在没有这种与受变化影响的状态和动作的“亲身经验”时适应环境变化。
模型变化会自动地（通过规划）改变策略。
规划可以确定环境变化的后果，即使这些变化从未在智能体自身经验中相互连接。
例如，仍以图 14.5 的迷宫任务为例，设想一只已有转移和奖励模型的大鼠被直接放进 $S_2$ 右侧目标箱，发现那里现在的奖励值从 4 变为 1。
即使找到该目标箱所需的迷宫动作选择并未涉及，大鼠的奖励模型也会改变。
规划过程会把新奖励知识用于跑迷宫，而不需要在迷宫中获得额外经验；在本例中，策略会改为在 $S_1$ 和 $S_3$ 都向右转，以获得回报 3。

正是这一逻辑构成了动物结果贬值实验的基础。
这类实验结果提供了关于动物是否学会习惯、还是其行为处于目标导向控制之下的线索。
结果贬值实验类似于潜伏学习实验，因为奖励会从一个阶段到下一个阶段发生变化。
在初始有奖励学习阶段之后，某个结果的奖励价值被改变，包括变为零甚至负值。

Adams 和 Dickinson 进行过这种类型的一个早期重要实验。
他们通过工具性条件化训练大鼠，直到大鼠在训练箱中积极按压杠杆以获得蔗糖颗粒。
然后把大鼠放入同一箱子，但收回杠杆，并允许非依赖性取食，也就是颗粒独立于大鼠动作而可获得。
在自由接触颗粒 15 分钟后，一组大鼠被注射会诱发恶心的氯化锂。
这重复了三次，在最后一次中，被注射的大鼠没有一只食用非依赖性颗粒，表明颗粒的奖励价值已经降低，即颗粒被贬值了。
第二天的下一阶段，大鼠再次被放入箱子并接受消退训练：反应杠杆重新出现，但与颗粒分配器断开，因此按压不会释放颗粒。
问题是，那些颗粒奖励价值被降低的大鼠，在没有经历“按压杠杆导致被贬值奖励”的情况下，是否会比未降低颗粒奖励价值的大鼠更少按压杠杆。
结果是，被注射大鼠从消退试次一开始，其反应率就显著低于未注射大鼠。

Adams 和 Dickinson 得出结论：被注射大鼠通过一张把按杠杆连接到颗粒、把颗粒连接到恶心的认知地图，把按杠杆与随后的恶心联系起来。
因此，在消退试次中，大鼠“知道”按杠杆的后果将是它们不想要的东西，于是一开始就减少按杠杆。
重要点在于，它们减少按杠杆，并不是因为曾经直接经历按杠杆后生病；它们生病时没有杠杆在场。
它们似乎能够把行为选择的结果知识（按杠杆会得到颗粒）与结果的奖励价值（颗粒应避免）结合起来，并据此改变行为。
并非所有心理学家都同意对这类实验的“认知”解释，这也不是解释这些结果的唯一可能方式，但基于模型规划的解释被广泛接受。

没有任何东西阻止智能体同时使用无模型和基于模型的算法，而且有充分理由同时使用二者。
我们从自身经验知道，经过足够重复，目标导向行为往往会变成习惯性行为。
实验表明，大鼠也会发生这种情况。
Adams 进行实验，考察延长训练是否会把目标导向行为转化为习惯性行为。
他通过比较接受不同训练量的大鼠在结果贬值后的反应来做到这一点。
如果延长训练使大鼠相比训练较少的大鼠对贬值不那么敏感，这就说明延长训练使行为更习惯化。
简化地说，一组大鼠训练到做出 100 次有奖励杠杆按压；另一组过度训练组训练到做出 500 次有奖励杠杆按压。
之后，两组大鼠的颗粒奖励价值都被降低。
然后两组都接受消退训练。
Adams 的问题是：贬值对过度训练大鼠杠杆按压率的影响，是否小于对非过度训练大鼠的影响。
结果是，贬值强烈降低了非过度训练大鼠的杠杆按压率；相反，对过度训练大鼠，贬值几乎没有影响，甚至似乎使按压更有力。
这一结果提示，非过度训练大鼠以目标导向方式行动，且对自己动作结果的知识敏感；而过度训练大鼠已经形成按杠杆的习惯。

从计算视角看待这一结果和其他类似结果，有助于理解为什么可以预期动物在某些情形下表现为习惯性，在另一些情形下表现为目标导向，并且为什么随着继续学习会从一种控制模式转向另一种。
虽然动物使用的算法无疑不完全匹配本书介绍的算法，但考虑各种强化学习算法所蕴含的权衡，可以帮助理解动物行为。
计算神经科学家 Daw、Niv 和 Dayan 提出的一个思想是，动物同时使用无模型和基于模型过程。
每个过程提出一个动作，而最终执行的动作由两个过程中被判断为更可信的过程决定；可信度由学习过程中维护的置信度量确定。
学习早期，基于模型系统的规划过程更可信，因为它把短期预测串接起来，而短期预测比无模型过程的长期预测能用较少经验变得准确。
但随着经验继续，无模型过程变得更可信，因为规划容易因模型不准确和使规划可行所需的捷径而出错，例如各种形式的“树剪枝”，即移除不太有希望的搜索树分支。
按这一思想，随着经验积累，应当预期从目标导向行为转向习惯性行为。
关于动物如何在目标导向控制与习惯控制之间仲裁，还提出过其他思想；行为和神经科学研究仍在继续考察这些及相关问题。

无模型和基于模型算法之间的区别正在证明对这类研究有用。
人们可以在抽象设定中考察这类算法的计算含义，揭示每种类型的基本优点和限制。
这既能提出也能澄清问题，从而指导实验设计，推动心理学家理解习惯性和目标导向的行为控制。

\subsection*{14.7\quad 小结}

本章的目标是讨论强化学习与心理学中动物学习实验研究之间的对应。
开头我们强调，本书所描述的强化学习并不旨在建模动物行为的细节。
它是一个抽象计算框架，从人工智能和工程视角考察理想化情形。
但是，许多基本强化学习算法受心理学理论启发；在某些情况下，这些算法也促进了新的动物学习模型的发展。
本章描述了其中最显著的对应。

强化学习中预测算法与控制算法的区别，平行于动物学习理论中经典（或巴甫洛夫）条件化与工具性条件化的区别。
工具性与经典条件化实验的关键差异是：前者中强化刺激依赖于动物行为，而后者中并不依赖。
通过 TD 算法学习预测对应于经典条件化；我们描述了经典条件化 TD 模型，这是强化学习原则解释动物学习行为某些细节的一个例子。
该模型把有影响力的 Rescorla--Wagner 模型推广到包含时间维度，使单个试次内事件影响学习，并解释二阶条件化，即强化刺激的预测物本身成为强化物。
它也是关于脑内多巴胺神经元活动的一种有影响力观点的基础；第 15 章将讨论这一点。

试错学习位于强化学习控制方面的基础处。
我们介绍了 Thorndike 对猫和其他动物的实验细节，这些实验导致了效果律；本章和第 1 章都讨论过该定律。
我们指出，在强化学习中，探索不必局限于“盲目摸索”；只要存在某种探索，试次可以由使用先天知识和先前学习知识的复杂方法生成。
我们讨论了 B. F. Skinner 称为塑造的训练方法，其中逐步改变奖励依赖关系，使动物逐次接近期望行为。
塑造不仅对动物训练不可或缺，也是训练强化学习智能体的有效工具。
这里还与动物动机状态这一思想相联系；动机状态影响动物会接近或避开什么，以及哪些事件对动物有奖赏性或惩罚性。

本书介绍的强化学习算法包含两个处理延迟强化问题的基本机制：资格迹和通过 TD 算法学习的价值函数。
两种机制在动物学习理论中都有先驱。
资格迹类似于早期理论中的刺激迹，价值函数则对应于次级强化在提供近乎即时的评价反馈方面的作用。

本章接着讨论的对应，是强化学习的环境模型与心理学家所谓认知地图之间的对应。
20 世纪中期进行的实验声称展示了动物学习认知地图的能力；认知地图可作为状态--动作联结的替代或补充，并在后来用于指导行为，尤其是在环境意外改变时。
强化学习中的环境模型类似于认知地图，因为它们可以由监督学习方法学得而不依赖奖励信号，随后用于规划行为。

强化学习在无模型与基于模型算法之间的区别，对应于心理学中习惯性行为与目标导向行为之间的区别。
无模型算法通过访问存储在策略或动作价值函数中的信息来作决策；基于模型的方法则使用智能体环境模型进行前瞻规划来选择动作。
结果贬值实验提供关于动物行为是习惯性的还是处于目标导向控制之下的信息。
强化学习理论帮助澄清了关于这些问题的思考。

动物学习显然启发了强化学习；但作为机器学习的一种，强化学习旨在设计和理解有效学习算法，而不是复制或解释动物行为细节。
我们聚焦于那些以清晰方式关联到预测和控制问题求解方法的动物学习方面，强调强化学习与心理学之间富有成果的双向思想流动，而没有深入许多占据动物学习研究者注意力的行为细节和争议。
未来强化学习理论与算法的发展，很可能会利用与动物学习更多特征之间的联系，因为这些特征的计算效用会得到更好认识。
我们预期，强化学习与心理学之间的思想流动会继续为两个学科带来成果。

强化学习与心理学及其他行为科学领域之间的许多联系超出了本章范围。
我们很大程度上省略了与决策心理学的联系；决策心理学关注学习发生之后动作如何被选择或决策如何作出。
我们也没有讨论行为的生态和进化方面，即动物如何彼此相处、如何与物理环境相处，以及它们的行为如何贡献于进化适应度。
优化、MDP 和动态规划在这些领域中都占重要位置；我们对智能体与动态环境交互的强调，也连接到复杂“生态”中智能体行为的研究。
本书省略的多智能体强化学习与行为的社会方面有联系。
尽管这里没有展开，强化学习绝不应被理解为忽视进化视角。
强化学习中没有任何东西暗示学习和行为是一张白纸。
事实上，工程应用经验已经突出表明，把类似于进化赋予动物的知识内建到强化学习系统中非常重要。

\section*{文献与历史备注}

Ludvig、Bellemare 和 Pearson，以及 Shah，分别从心理学和神经科学语境综述了强化学习。
这些文献是本章以及下一章强化学习与神经科学讨论的有用补充。

\noindent\textbf{14.1}
关于经典条件化与工具性条件化的相互作用，特别是经典条件化反应与工具性反应冲突的情形，Dayan、Niv、Seymour 和 Daw 作了重点研究。
他们提出了一个 Q 学习框架来建模这种相互作用的若干方面。
Modayil 和 Sutton 使用移动机器人展示了把固定反应与在线预测学习结合起来的控制方法的有效性；他们称之为巴甫洛夫控制，并强调它不同于强化学习通常的控制方法，因为它基于预测性地执行固定反应，而不是奖励最大化。

\noindent\textbf{14.2.1}
Kamin 最先报告了经典条件化中的阻滞，现在通常称为 Kamin 阻滞。
Moore 和 Schmajuk 对阻滞现象、它激发的研究及其对动物学习理论的持久影响给出了很好的概述。
关于兔瞬膜反应的二阶条件化及其与串行复合刺激条件化的关系，也有专门研究。

\noindent\textbf{14.2.2}
Rescorla--Wagner 模型中“动物在感到意外时才学习”的思想来自 Kamin。
除 Rescorla 和 Wagner 的模型外，经典条件化模型还包括 Klopf、Grossberg、Mackintosh、Moore 和 Stickney、Pearce 和 Hall、Courville、Daw 和 Touretzky 等人的模型。
Schmajuk 综述了经典条件化模型，Wagner 给出了关于 Rescorla--Wagner 模型及类似元素式学习理论的现代心理学视角。

\noindent\textbf{14.2.3}
经典条件化 TD 模型的早期版本出现在 Sutton 和 Barto 的工作中，其中也包含“时间优先性压倒阻滞”的早期预测，后来由 Kehoe、Schreurs 和 Graham 在兔瞬膜制备中证实。
这些早期工作还最早认识到 Rescorla--Wagner 模型与最小均方（LMS）或 Widrow--Hoff 学习规则几乎相同。
该早期模型在 Sutton 发展 TD 算法之后被修订，并最终以 TD 模型形式更完整呈现。
Moore 及其同事还研究了 TD 模型及其可能的神经实现。
Klopf 的驱动--强化经典条件化理论扩展了 TD 模型，以处理更多实验细节，例如获得曲线的 S 形。

\noindent\textbf{14.2.4}
Ludvig、Sutton 和 Kehoe 评价了 TD 模型在此前未探索的经典条件化任务中的表现，并考察了各种刺激表征的影响，其中包括他们较早提出的微刺激表征。
关于各种刺激表征及其可能神经实现对反应时序和形态影响的更早研究，可见 Moore 及其同事的工作。
第 353--355 页的图改编自 Sutton 和 Barto 的早期论文。

\noindent\textbf{14.3}
第 1.7 节包含关于试错学习和效果律历史的说明。
强化学习中的塑造示例包括杆平衡任务和若干机器人学习任务。
Ng 及其同事也使用“塑造”一词，但其意义与 Skinner 稍有不同，重点在于如何改变奖励信号而不改变最优策略集合。
关于学习和动机相互作用的复杂性，Dickinson 和 Balleine 有讨论；Wise、Daw、Shohamy、McClure、Niv、Dayan 等人的工作把动机与强化学习理论的若干方面联系起来，并提出了与强化学习框架相关的行为活力理论。

\noindent\textbf{14.4}
Hull 在耶鲁的学生和合作者 Spence 进一步阐述了高阶强化在处理延迟强化问题中的作用。
非常长延迟下的学习，例如延迟可达数小时的味觉厌恶条件化，引出了干扰理论作为衰减迹理论的替代。
关于延迟强化下学习的其他观点，还引入了意识和工作记忆的作用。

\noindent\textbf{14.5}
Thistlethwaite 对其出版时为止的潜伏学习实验作了广泛综述。
Ljung 概述了工程中的模型学习或系统辨识技术。
Gopnik 等人提出了儿童如何学习模型的贝叶斯理论。

\noindent\textbf{14.6}
习惯性和目标导向行为与无模型和基于模型强化学习之间的联系，最早由 Daw、Niv 和 Dayan 提出。
用于解释习惯性和目标导向行为控制的假想迷宫任务，基于 Niv、Joel 和 Dayan 的说明。
Dolan 和 Dayan 综述了与该问题相关的四代实验研究，并讨论如何基于强化学习的无模型/基于模型区分继续推进。
Dickinson、Balleine 等人讨论了与这种区分相关的实验证据；也有研究者主张无模型过程也能解释结果贬值实验。

\noindent\textbf{关于术语的说明}
心理学中“强化”的传统含义，是动物在与另一个刺激或反应具有适当时间关系时接收到某个刺激（或经历某个刺激的省略）后，某种行为模式被增强，即强度或频率增加。
强化产生会保留在未来行为中的变化。
心理学中有时也把“强化”用于产生行为持久变化的过程，无论变化是增强还是削弱行为模式。
让强化也指削弱，虽与日常语言和心理学传统用法不完全一致，但这是一个有用扩展，本书也采用了它。
无论哪种情况，被认为导致行为变化的刺激称为强化物。

心理学家通常不像我们这样使用“强化学习”这一特定短语。
动物学习先驱可能把强化和学习视为同义，因此同时使用二者会显得冗余。
我们的用法沿袭计算和工程研究中的用法，主要受 Minsky 影响。
不过，这一短语近来在心理学和神经科学中也越来越常见，很可能是因为强化学习算法与动物学习之间出现了强平行关系，也就是本章和下一章描述的这些平行关系。

按通常用法，奖励是动物会接近并愿意为之工作的对象或事件。
奖励可以因动物“好”的行为而给出，也可以为了让动物行为“更好”而给出。
类似地，惩罚是动物通常会避免的对象或事件，它作为“坏”行为的结果给出，通常为了改变该行为。
初级奖励来自进化内建于动物神经系统的机制，以提高其生存和繁殖机会；例如营养食物的味道、性接触、成功逃脱，以及祖先历史中预测繁殖成功的许多其他刺激和事件。
高阶奖励由预测初级奖励的刺激提供，直接预测或间接通过预测其他预测初级奖励的刺激。
如果一个刺激的奖赏性来自直接预测初级奖励，则它是次级奖励。

本书把 $R_t$ 称为“时刻 $t$ 的奖励信号”，有时简称“时刻 $t$ 的奖励”，但我们不把它看成智能体环境中的对象或事件。
因为 $R_t$ 是一个数，而不是对象或事件，所以它更像神经科学中的奖励信号；后者是脑内部的信号，例如神经元活动，会影响决策和学习。
这个信号可能在动物感知到有吸引力或厌恶性的对象时被触发，也可能由外部环境中并不存在的东西触发，例如记忆、想法或幻觉。
由于我们的 $R_t$ 可以为正、为负或为零，把负的 $R_t$ 称为惩罚、把零称为中性信号也许更好；但为简单起见，我们通常避免这些术语。

在强化学习中，产生所有 $R_t$ 的过程定义了智能体试图解决的问题。
智能体目标是在时间上尽量保持 $R_t$ 的大小较大。
在这一点上，如果把动物面对的问题理解为在其一生中获得尽可能多的初级奖励，那么 $R_t$ 类似于动物的初级奖励。
不过，如第 15 章将提示的，动物脑中不太可能存在一个像 $R_t$ 这样的单一“总控”奖励信号。

并非所有强化物都是奖励或惩罚。
有时强化并不是因为动物接收到一个通过把行为标记为好或坏来评价行为的刺激。
某种行为模式可以被一个无论动物如何行为都会到来的刺激强化。
如第 14.1 节所述，强化物的给出是否依赖先前行为，是工具性或操作性条件化实验与经典或巴甫洛夫条件化实验之间的定义性差异。
两类实验中都有强化起作用，但只有前者中的强化才是评价过去行为的反馈。

当下一章讨论这些信号的神经相关物时，奖励信号与强化信号之间的区别非常关键。
对我们来说，某一具体时刻的强化信号与奖励信号一样，是一个正数、负数或零。
强化信号是指导学习算法改变智能体策略、价值估计或环境模型的主要因素。
在我们看来，最有意义的定义是：某一时刻的强化信号是一个数，它乘以（可能连同某些常数）一个向量，从而决定某个学习算法中的参数更新。

对某些算法来说，奖励信号本身就是参数更新方程中的关键乘子；这些算法的强化信号与奖励信号相同。
但对本书讨论的大多数算法来说，强化信号除奖励信号外还包括其他项。
例如，TD 状态价值学习的强化信号是 TD 误差
\[
\delta_t=R_{t+1}+\gamma V(S_{t+1})-V(S_t).
\]
在这个强化信号中，$R_{t+1}$ 是初级强化贡献，而预测价值的时间差分 $V(S_{t+1})-V(S_t)$（或动作价值的类似时间差分）是条件强化贡献。
因此，当 $V(S_{t+1})-V(S_t)=0$ 时，$\delta_t$ 表示“纯”初级强化；当 $R_{t+1}=0$ 时，它表示“纯”条件强化；但它常常表示二者的混合。

一个可能的混淆来源是著名心理学家 B. F. Skinner 及其追随者使用的术语。
对 Skinner 来说，正强化发生在动物行为后果增加该行为频率时；惩罚发生在行为后果降低该行为频率时。
负强化发生在行为导致厌恶刺激移除，从而增加该行为频率时。
另一方面，负惩罚发生在行为导致食欲性刺激移除，从而降低该行为频率时。
我们并不特别需要这些区分，因为我们的方法比这更抽象，奖励和强化信号都允许取正值和负值。
尤其要注意，当我们的强化信号为负时，它并不等同于 Skinner 的负强化。

另一方面，人们也常指出，用单一数字依其符号作为奖励或惩罚信号，与动物的食欲系统和厌恶系统具有质性不同性质、并涉及不同脑机制这一事实并不一致。
这指出了强化学习框架未来可能发展的方向：利用分离的食欲系统和厌恶系统的计算优势。
但目前我们暂且略过这些可能性。

另一个术语差异是我们如何使用“动作”一词。
对许多认知科学家来说，动作是有目的的，意思是它来自动物关于相关行为与该行为后果之间关系的知识。
动作是目标导向的，是决策的结果；相对地，反应由刺激触发，是反射或习惯的结果。
我们使用“动作”一词时，并不区分别人所谓动作、决策和反应。
这些是重要区分，但对我们来说，它们被无模型与基于模型强化学习算法之间的差异所涵盖；上文在第 14.6 节讨论习惯性和目标导向行为时已经说明。

本书中经常使用的另一个词是“控制”。
我们所说的控制，与动物学习心理学家所说的控制完全不同。
我们说控制，是指智能体影响其环境，以带来智能体偏好的状态或事件；智能体对其环境施加控制。
这是控制工程师使用的控制含义。
相反，在心理学中，控制通常意味着动物行为受到它接收的刺激（刺激控制）或它经历的强化时程的影响，也就是被这些东西控制。
这里是环境控制智能体。
当然，当智能体与环境交互时，这两个方向的控制都在发挥作用；但我们的焦点是作为控制者的智能体，而不是作为控制者的环境。
一种等价且也许更有启发性的观点是，智能体实际上在控制它从环境接收的输入。
这并不是心理学家所说的刺激控制。

有时，人们把强化学习理解为仅指不涉及价值函数或环境模型、直接从奖励（和惩罚）学习策略。
这就是心理学家所谓刺激--反应学习，或 S--R 学习。
但对我们以及今天大多数心理学家来说，强化学习比这宽泛得多；除了 S--R 学习，它还包括涉及价值函数、环境模型、规划以及其他通常被认为属于心理功能中更认知一侧的过程的方法。

\end{document}
